\section{Related Work}
\label{sec:related}

\textbf{MCP Server Security and Risk Assessment.}
Prior work surveys security risks across the MCP
ecosystem~\cite{hou2026mcp} and audits open-source MCP servers for
software weaknesses, MCP-specific vulnerabilities, and maintainability
problems~\cite{hasan2025firstglance}. Several studies examine attacks in which
a malicious server, tool description, or external input influences the
agent~\cite{invariant2025toolpoisoning,shi2025toolhijacker,
zhao2025parasites}. We study the opposite direction: an agent sends a call that
may harm the MCP server or an organizational resource reached through it.
Server-to-agent attacks are outside our scope. We assume only that each tool
performs the operation represented by its declared interface and resource
mapping.

Other work examines the capabilities and weaknesses of MCP servers. Li et al.
analyze 2,562 MCP applications and measure their use of file, network, system,
and memory APIs~\cite{li2025privilege}. Huang et al. combine static pattern
matching with sandboxed execution to detect capabilities such as file access,
network requests, and command execution~\cite{huang2026auditing}. MCP-in-SoS
scans 222 server repositories for code weaknesses, maps them to CWE and CAPEC
entries, and derives repository-level risk from the likelihood and severity
those catalogs assign~\cite{kumar2026mcpinsos}. All three assess the server as
an artifact, independently of the resources a given installation holds, so the
same server receives the same assessment wherever it is deployed.

ConLeash checks whether an MCP call stays within an approved consent
boundary~\cite{li2026conleash}. However, it does not estimate the risk or
possible impact of a call that is already allowed. Our work addresses this gap
by assessing the risk of each invocation from the server side.

\textbf{Defenses for MCP and tool-using agents.}
IGAC checks whether an MCP invocation and its arguments remain within the
user's stated intent, while MTGuard compares invocation parameters, runtime
behavior, and results against safety policies~\cite{zhu2026igac,mtguard2026}.
Progent and MiniScope enforce least-privilege authorization, whereas
AgentGuardian learns admissible tool inputs and invocation sequences from
normal executions~\cite{shi2025progent,zhu2025miniscope,abaev2026agentguardian}.
These systems check whether a tool call follows the user's request, the
security policy, or the agent's normal behavior. However, a call can pass all
these checks and still cause serious damage. For example, an agent may be
allowed to delete files, but deleting a temporary file and deleting a critical
production file do not have the same risk; our work measures this difference at
the MCP server.

\textbf{Access control.}
RBAC assigns permissions through roles, while ABAC evaluates attributes of the
caller, resource, requested action, and environment
conditions~\cite{sandhu1996rbac,hu2014abac}. These models can express
fine-grained policies, but they do not directly estimate the possible
consequence of a permitted request.

Risk-adaptive access control adds a calculated risk value to the authorization
decision~\cite{cheng2007fuzzymls}. Kandala et al. place risk within an
attribute-based framework, but leave the concrete risk-calculation method
outside their scope~\cite{kandala2011radac}. Ni et al. use fuzzy rules over
subject and object security information to estimate the risk of an access
operation~\cite{ni2010fuzzy}. Shaikh et al. dynamically calculate trust and
risk for each subject--object pair from the user's past
behavior~\cite{shaikh2012dynamic}. Atlam et al. survey the wider family of
risk-based access-control models~\cite{atlam2020riskbased}.

For agentic systems, Fleming et al. use an LLM judge to create a task policy,
combine static risk values for its tools and resources, and estimate
uncertainty~\cite{fleming2025tbac}. A limitation is that risk is computed for
the whole task and is not updated for each call. For example, reading one
database record and exporting a full table may use the same tool risk value.
Our work instead scores each MCP invocation using its runtime details.

\textbf{Risk scoring for MCP and agents.}
CVSS scores the severity of software vulnerabilities. A study of 196 evaluators
applying CVSS v3.1 found differences between evaluators, and a follow-up found
that 68\% assigned a different severity level to vulnerabilities they had
scored before~\cite{first2023cvss,wunder2024cvss}. AIVSS extends vulnerability
scoring to agentic AI systems~\cite{owasp2026aivss}. These frameworks score
vulnerabilities rather than concrete tool calls.

FIPS 199 and NIST SP 800-60 classify information types and systems by the
potential impact of losing confidentiality, integrity, or
availability~\cite{nist2004fips199,nist2008sp80060}. These classifications
provide an organizational view of consequence that our method uses when
assessing the resource affected by a call.

Work closer to our setting studies risk at different stages. MCP-in-SoS scores
weaknesses in an MCP server repository~\cite{kumar2026mcpinsos}, while
AgenTRIM assigns risk levels to tools and applies stronger checks to riskier
ones~\cite{betser2026agentrim}. STARS scores proposed skill invocations using
the request, capability, and runtime context before
execution~\cite{zhang2026stars}. Owiredu-Ashley instead grades the harm of
actions recorded after execution~\cite{owireduashley2026severity}. Together,
these methods answer different risk questions: they assess server weaknesses,
tool capabilities, invocation safety, or harm after execution. They do not
provide a common measure of the consequence to the organization affected by
the call. For example, deleting a temporary cache and deleting payroll records
may use the same tool and both be irreversible, but their organizational impact
is very different. Our work connects each MCP call to the organization's
resource policy and estimates its risk before execution.
